% COLLECTION_METADATA: {"schema_version": 1, "kind": "reused_writeup", "independently_reviewed": false, "primary_source": "attacks/erdos/gpt_pro_5.2/89.tex", "source_paths": ["attacks/erdos/gpt_5.2/89.tex", "attacks/erdos/codex_5.2_extra_high/89.tex", "attacks/erdos/gpt_pro_5.2/89.tex"], "source_model": "gpt_pro_5.2", "source_completion": null, "source_claim": "unresolved"}
\section*{Erd\H{o}s Problem 89 --- collection record}

\subsection*{PROVENANCE AND STATUS}
Official problem: https://www.erdosproblems.com/89

Saved collaborative-database status: open.
Snapshot checked: 2026-09-23T17:22:24Z.
Latest status and formalization links: https://teorth.github.io/erdosproblems/

This is an attributed collection of existing writeups. The model-folder names below identify their original attribution; placement in GPT\_6\_Astra\_Ultra does not attribute their mathematical work to that model. All locally available source versions selected for this problem are reproduced below, without editing their text.

Proof claims, computations, literature statements, and completion estimates in the reused text have not been independently verified in this collection pass. A source claiming a solution does not change the saved upstream status.

Primary source (latest numbered version in the preferred model folder): \texttt{attacks/erdos/gpt\_pro\_5.2/89.tex}.

\subsection*{REUSED SOURCE 1}
Original attribution: \texttt{gpt\_5.2}.
Original file: \texttt{attacks/erdos/gpt\_5.2/89.tex}.

% BEGIN REUSED SOURCE: attacks/erdos/gpt_5.2/89.tex

OUTPUT
1) FORMAL RESTATEMENT
Definitions/Conventions.
Let $A\subset\mathbb R^2$ be a set of $n$ distinct points.
Let
\[D(A):=\big|\{|x-y|: x,y\in A,\ x\ne y\}\big|\]
be the number of distinct Euclidean distances determined by unordered pairs of points of $A$.

Question.
Is it true that for all $n$ and all $|A|=n$ one has
\[D(A)\ \gg\ \frac{n}{\sqrt{\log n}}?\]

2) QUICK LITERATURE/CONTEXT CHECK
Web browsing is not available in this session.
From the statement only:
- Guth--Katz proved $D(A)\gg n/\log n$.
I do not re-prove that deep result here.

3) ATTACK PLAN
Proof track ideas:
1. Use the standard energy/Cauchy--Schwarz reduction: distinct distances are large if no distance occurs too many times.
2. Bound the number of repeated distances via incidence geometry (point-circle incidences); the full problem uses sophisticated polynomial methods.

Disproof track ideas:
1. Try to build configurations with very few distinct distances; the integer grid is the classical near-extremizer.

Chosen path in this attempt: provide a fully proved warm-up bound $D(A)\gg \sqrt n$ using a K\H{o}v\'ari--S\'os--Tur\'an argument for repeated distances at a fixed radius.

4) WORK
Lemma 1 (K\H{o}v\'ari--S\'os--Tur\'an for $K_{2,3}$-free bipartite graphs; a usable form).
Let $B$ be a bipartite graph with left part $L$ and right part $R$, where $|L|=|R|=n$.
Assume $B$ contains no copy of $K_{2,3}$ with the two-vertex side in $L$ and the three-vertex side in $R$ (equivalently, any two vertices in $L$ have at most $2$ common neighbors in $R$).
Then
\[e(B)\le 2n^{3/2}+2n.\]

Proof.
Let the degrees of left vertices be $d(\ell)$ for $\ell\in L$.
Count length-2 paths $\ell-r-\ell'$ with $\ell\ne \ell'$ and $r\in R$.
For a fixed $r\in R$, the number of unordered pairs of distinct neighbors of $r$ in $L$ is $\binom{d(r)}{2}$, so the total number of such 2-paths equals
\[P:=\sum_{r\in R}\binom{d(r)}{2}.\tag{1}\]
On the other hand, for each unordered pair $\{\ell,\ell'\}\subseteq L$, the number of common neighbors in $R$ is at most $2$ by the $K_{2,3}$-free assumption.
Thus each pair $\{\ell,\ell'\}$ contributes at most $2$ to $P$, so
\[P\le 2\binom{n}{2} < n^2.\tag{2}\]

Now relate $P$ to the total edge count $e=e(B)$.
By convexity (or Cauchy--Schwarz), $\sum_{r\in R} d(r)^2 \ge e^2/n$ since $\sum_{r} d(r)=e$.
Also $\binom{d}{2}=\frac{d^2-d}{2}$, so
\[P=\sum_{r}\binom{d(r)}{2}=\frac12\sum_r d(r)^2 - \frac12\sum_r d(r)\ge \frac12\cdot \frac{e^2}{n}-\frac12 e.\tag{3}\]
Combine (2) and (3):
\[\frac{e^2}{2n}-\frac{e}{2} \le n^2.\]
Multiply by $2n$:
\[e^2 - en \le 2n^3.\]
Complete the square: $(e-\tfrac{n}{2})^2 \le 2n^3 + \tfrac{n^2}{4} \le (2n^{3/2}+n)^2$ for $n\ge 1$.
Thus $e-\tfrac{n}{2}\le 2n^{3/2}+n$, i.e. $e\le 2n^{3/2}+\tfrac{3n}{2}<2n^{3/2}+2n$.
\qed

Lemma 2 (Warm-up distinct-distance bound $D(A)\gg \sqrt n$).
For any set $A$ of $n$ points in $\mathbb R^2$,
\[D(A)\ \ge\ c\,\sqrt n\]
for an absolute constant $c>0$.

Proof.
For a distance value $r>0$, let $m_r$ be the number of unordered pairs $\{x,y\}\subset A$ with $|x-y|=r$.
Then
\[\sum_r m_r = \binom{n}{2}.\tag{4}\]
It suffices to show that every $m_r\ll n^{3/2}$, because then
\[\binom{n}{2}=\sum_r m_r \le D(A)\cdot \max_r m_r \ll D(A)\,n^{3/2},\]
giving $D(A)\gg n^2/n^{3/2}=\sqrt n$.

Fix $r>0$. Build a bipartite graph $B_r$ with left part $L=A$ (centers) and right part $R=A$ (points), and draw an edge from $x\in L$ to $y\in R$ iff $|x-y|=r$.
Then $e(B_r)=2m_r$ (each unordered pair contributes two directed edges).

Key geometric fact: any two circles in the plane intersect in at most $2$ points.
Thus for two distinct centers $x\ne x'$, the circles of radius $r$ around $x$ and $x'$ have at most two intersection points, meaning $x$ and $x'$ have at most $2$ common neighbors in $R$.
Equivalently, $B_r$ is $K_{2,3}$-free in the sense of Lemma 1.
Therefore Lemma 1 gives
\[2m_r=e(B_r)\le 2n^{3/2}+2n,\]
so $m_r\le n^{3/2}+n$ for every $r$.
Substitute into (4) to conclude $D(A)\ge \binom{n}{2}/(n^{3/2}+n)\gg \sqrt n$. \qed

FAST REALITY CHECK
The desired bound is $D(A)\gg n/\sqrt{\log n}$, much stronger than the $\Omega(\sqrt n)$ bound proved here.
The $\Omega(\sqrt n)$ bound matches what one would get from the ``easy'' $O(n^{3/2})$ unit-distance bound; improving beyond that requires deep incidence geometry.

5) VERIFICATION
- Lemma 1: the only structural assumption used is ``any two left vertices have at most 2 common neighbors,'' which is exactly the $K_{2,3}$-free condition.
- Lemma 2: the geometric input ``two circles meet in at most 2 points'' is correct, and the translation to the bipartite graph $B_r$ is explicit.

6) FINAL
**UNRESOLVED**

(i) Strongest fully proved partial result:
I proved the warm-up bound $D(A)\gg \sqrt n$ for all $n$-point sets in the plane (Lemma 2).

(ii) First gap:
Bound the number of repeated distances more sharply than $O(n^{3/2})$ per distance, or control the \emph{sum} of squared multiplicities $\sum_r m_r^2$ well enough to upgrade Cauchy--Schwarz to $n/\sqrt{\log n}$.

(iii) Top 3 next moves (concrete):
1. Use the Cauchy--Schwarz framework $D(A)\ge \binom{n}{2}^2 / \sum_r m_r^2$ and attempt to bound $\sum_r m_r^2$ via incidence bounds between points and circles.
2. Specialize to structured candidate extremizers (integer grids) and understand their distance multiplicities; attempt to prove the $n/\sqrt{\log n}$ upper behavior for grids (number theory of sums of two squares).
3. Explore whether stronger forbidden-subgraph conditions (e.g. no three collinear, no four concyclic) allow sharper incidence bounds as a stepping stone (cf. Problem \#98).

(iv) What a minimal counterexample would likely look like (structure/parameters):
A family of point sets $A_n$ with $|A_n|=n$ such that the number of distinct distances $D(A_n)$ is $o\!\left(n/\sqrt{\log n}\right)$,
meaning that many distances are repeated with unusually high multiplicity across $\Theta(n^2)$ pairs.


% END REUSED SOURCE: attacks/erdos/gpt_5.2/89.tex

\subsection*{REUSED SOURCE 2}
Original attribution: \texttt{codex\_5.2\_extra\_high}.
Original file: \texttt{attacks/erdos/codex\_5.2\_extra\_high/89.tex}.

% BEGIN REUSED SOURCE: attacks/erdos/codex_5.2_extra_high/89.tex

OUTPUT
1) FORMAL RESTATEMENT
Let $A \subset \mathbb{R}^2$ be a finite set of $n$ distinct points, $n \ge 2$.
Define the distinct distance set
  $D(A) := \{ \|x-y\| : x,y \in A,\ x\neq y\}$,
and let $d(A) := |D(A)|$.
We interpret $f(n) \gg g(n)$ as: there exist constants $c>0$ and $n_0$ such that
for all $n \ge n_0$, $f(n) \ge c\, g(n)$. The logarithm is natural; the base is irrelevant.

Claim (main): There exist $c>0$ and $n_0$ such that for all $n \ge n_0$ and all
$A \subset \mathbb{R}^2$ with $|A|=n$,
  $d(A) \ge c\, n/\sqrt{\log n}$.

Stronger variants (from statement) include pointwise/average versions:
Define $d(x) := |\{\|x-y\|: y\in A,\ y\neq x\}|$. One conjecture is
  $\sum_{x\in A} d(x) \gg n^2/\sqrt{\log n}$.

2) QUICK LITERATURE/CONTEXT CHECK
Browsing is not available. The problem statement itself asserts:
- A $\sqrt{n}\times\sqrt{n}$ integer grid shows the bound $n/\sqrt{\log n}$ is
  (up to constants) best possible.
- Guth--Katz (2015) proved the lower bound $d(A) \gg n/\log n$.
I do not re-derive those results here; I treat them only as context from the prompt.

3) ATTACK PLAN
Proof track:
1. Incidence geometry / polynomial method: analyze point-line incidences of
   perpendicular bisectors; known to yield Guth--Katz-type bounds.
2. Energy/graph methods: relate repeated distances to many isosceles triangles,
   then bound via incidence theorems.
3. Sum-product or additive combinatorics on distance sets for structured point
   sets (grids, Cartesian products).

Disproof track:
1. Construct point sets with unusually few distances; the integer grid is the
   canonical obstruction.
2. Examine lattice-like sets with many repeated distances (sum of two squares).
3. Attempt thin annulus configurations or product sets to push below
   $n/\sqrt{\log n}$ (likely impossible if conjecture is true).

Chosen path: With current tools, neither a full proof nor counterexample is
reachable here. I focus on rigorous small-case analysis and a concrete partial
lemma, and then record the first true gap.

4) WORK
Lemma 1 (Equidistant sets in the plane are small).
If $A \subset \mathbb{R}^2$ has the property that all pairwise distances are equal,
then $|A|\le 3$.

Proof.
Assume all pairwise distances in $A$ equal $r>0$. Pick distinct points $a,b,c \in A$.
Then $\triangle abc$ is equilateral with side length $r$.
Any point $x \in A$ must satisfy $\|x-a\|=\|x-b\|=r$, hence $x$ lies on the intersection
of the two circles centered at $a$ and $b$ with radius $r$. These two circles intersect
in at most two points: $c$ and its reflection $c'$ across the line $ab$.
If $x=c$, it is not new. If $x=c'$, then $\|x-c\| = \|c'-c\| = r\sqrt{3} \ne r$,
contradicting the equidistance assumption. Hence no fourth point exists.
Therefore $|A| \le 3$. ∎

Corollary 2 (At least two distances for $n\ge 4$).
If $|A|\ge 4$, then $d(A)\ge 2$.

Proof.
If $d(A)=1$, then all pairwise distances are equal, so by Lemma 1 we have
$|A|\le 3$, contradiction. ∎

This is a very weak bound but it is fully rigorous and does not rely on any
deep theorems.

5) VERIFICATION
Edge cases:
- $n=2$: $d(A)=1$ is possible (two points).
- $n=3$: $d(A)=1$ is possible (equilateral triangle).
- $n=4$: Corollary 2 forces $d(A)\ge 2$, consistent with examples.

Quantifiers:
Lemma 1 assumes all pairwise distances equal and uses circle intersection facts.
The circle-intersection argument is standard Euclidean geometry and does not
depend on any hidden assumptions.

No hidden axioms or external results were used beyond elementary geometry.

6) FINAL
**UNRESOLVED**

(i) Strongest fully proved partial result:
For $n\ge 4$, any set of $n$ points in $\mathbb{R}^2$ determines at least
two distinct distances (Corollary 2).

(ii) First gap:
Proving any super-constant lower bound on $d(A)$ (e.g., $d(A)\ge c\, n^\alpha$
with $\alpha>0$) for all $A$ requires nontrivial incidence geometry beyond the
elementary tools used here.

(iii) Top 3 next moves:
1. Formalize the isosceles triangle count bound and connect it to
   point-line incidence theorems.
2. Import or reproduce the Guth--Katz framework to recover $d(A)\gg n/\log n$.
3. Quantify the grid obstruction via sum-of-two-squares counting to show
   $d(A)\ll n/\sqrt{\log n}$ for the $\sqrt{n}\times\sqrt{n}$ grid.

(iv) Minimal counterexample structure (if the conjecture were false):
A sequence of point sets with $|A|=n$ and $d(A)=o(n/\sqrt{\log n})$, likely
highly structured (Cartesian product / lattice-like) and exhibiting many
repeated distances beyond what is possible in grids.

% END REUSED SOURCE: attacks/erdos/codex_5.2_extra_high/89.tex

\subsection*{REUSED SOURCE 3}
Original attribution: \texttt{gpt\_pro\_5.2}.
Original file: \texttt{attacks/erdos/gpt\_pro\_5.2/89.tex}.

% BEGIN REUSED SOURCE: attacks/erdos/gpt_pro_5.2/89.tex

\noindent\textbf{FORMAL RESTATEMENT.}
Let $A\subset\mathbb{R}^2$ be a set of $n\ge 2$ distinct points. Let
\[
D(A):=\bigl|\{\,|x-y|:\ x,y\in A,\ x\ne y\,\}\bigr|
\]
be the number of distinct Euclidean distances determined by unordered pairs from $A$.
Question: does there exist an absolute constant $c>0$ such that for all sufficiently large $n$ and all such $A$,
\[
D(A)\ge c\,\frac{n}{\sqrt{\log n}}\ ?
\]

\medskip
\noindent\textbf{QUICK LITERATURE/CONTEXT CHECK.}
The problem statement notes that a $\sqrt n\times \sqrt n$ integer grid would make the bound $\gg n/\sqrt{\log n}$ best possible, and that Guth--Katz proved the lower bound $\gg n/\log n$. I do not use external results beyond elementary geometry/graph counting.

\medskip
\noindent\textbf{ATTACK PLAN.}
\begin{itemize}
\item \emph{Proof track:} prove a weaker lower bound for $D(A)$ by bounding how many pairs can realise the \emph{same} distance, and then pigeonholing over distances.
\item \emph{Disproof track:} look for constructions with $o(n/\sqrt{\log n})$ distinct distances; the known grid heuristic suggests this is unlikely.
\end{itemize}
We carry out the weaker proof track to obtain $D(A)=\Omega(\sqrt n)$.

\medskip
\noindent\textbf{WORK.}

\smallskip
\noindent\textbf{Lemma 89.1 (few pairs at any fixed distance).}
Fix $r>0$. Let $A\subset\mathbb{R}^2$ with $|A|=n$.
Let $U_r(A)$ be the number of unordered pairs $\{x,y\}\subset A$ with $|x-y|=r$.
Then
\[
U_r(A)\le \frac{\sqrt 2}{2}\,n^{3/2}+\frac14 n.
\]
\textit{Proof.}
Consider the bipartite graph $G=(L\sqcup R,E)$ where $L$ is a copy of $A$ (viewed as \emph{circle centers}) and $R$ is another copy of $A$ (viewed as \emph{points}). Put an edge between $\ell\in L$ (center $x\in A$) and $r\in R$ (point $y\in A$) iff $|x-y|=r$. Then $|L|=|R|=n$ and $|E|$ counts \emph{directed} distance-$r$ pairs, so
\[
|E|=2U_r(A).
\]

\emph{Claim 1:} $G$ is $K_{2,3}$-free.
Indeed, take two distinct centers $x_1\ne x_2$ (vertices in $L$). The set of points $y\in A$ that are adjacent to both $x_1$ and $x_2$ are exactly the intersection points of the two circles of radius $r$ centered at $x_1$ and $x_2$. Two distinct circles in the plane intersect in at most two points: subtracting their equations yields a line, which meets a circle in at most two points. Hence $x_1,x_2$ have at most two common neighbors in $R$, so there cannot be three points adjacent to both, i.e. no $K_{2,3}$.

\emph{Claim 2 (a concrete Kővári--Sós--Turán bound for $K_{2,3}$-free graphs):}
If a bipartite graph has parts of size $n$ and is $K_{2,3}$-free, then
\[
|E|\le \sqrt 2\,n^{3/2}+\frac12 n.
\]
\emph{Proof of Claim 2.}
Let the degrees of vertices on the right side be $d_1,\dots,d_n$ so that $\sum_{i=1}^n d_i=|E|$.
For a right vertex with degree $d_i$, the number of unordered pairs of its left neighbors is $\binom{d_i}{2}$.
Summing over $R$ counts, for each unordered left pair $\{u,v\}$, the number of common right neighbors.
Since the graph is $K_{2,3}$-free, any left pair has at most $2$ common right neighbors, so
\[
\sum_{i=1}^n \binom{d_i}{2} \le 2\binom{n}{2}=n(n-1).
\]
Rewrite the left-hand side as
\[
\sum_{i=1}^n \binom{d_i}{2}=
\frac12\sum_{i=1}^n(d_i^2-d_i)
=
\frac12\Bigl(\sum_{i=1}^n d_i^2\Bigr)-\frac12|E|.
\]
By Cauchy--Schwarz,
\[
\sum_{i=1}^n d_i^2 \ge \frac1n\Bigl(\sum_{i=1}^n d_i\Bigr)^2 = \frac{|E|^2}{n}.
\]
Therefore
\[
\frac12\Bigl(\frac{|E|^2}{n}\Bigr)-\frac12|E|
\le n(n-1)
\le n^2,
\]
which implies $|E|^2 - n|E| -2n^3\le 0$.
Solving this quadratic inequality in $|E|$ gives
\[
|E|\le \frac{n+\sqrt{n^2+8n^3}}{2}\le \frac{n+\sqrt{8}\,n^{3/2}}{2}=\sqrt 2\,n^{3/2}+\frac12 n.
\]
This proves Claim 2.

Combining Claim 1 and Claim 2 yields $|E|\le \sqrt 2\,n^{3/2}+\tfrac12 n$, hence
\[
U_r(A)=\frac{|E|}{2}\le \frac{\sqrt 2}{2}\,n^{3/2}+\frac14 n.
\]
\hfill $\square$

\smallskip
\noindent\textbf{Lemma 89.2 (a weak distinct-distance lower bound).}
For every finite $A\subset\mathbb{R}^2$ with $|A|=n\ge 2$,
\[
D(A)\ge \frac{n-1}{\sqrt 2\,\sqrt n+\tfrac12}\,\sqrt n \ge \frac{1}{2\sqrt 2}\,\sqrt n\quad\text{for }n\ge 4.
\]
In particular $D(A)=\Omega(\sqrt n)$.
\textit{Proof.}
Let the distinct distance values determined by $A$ be $r_1,\dots,r_m$ where $m=D(A)$.
Then every unordered pair of distinct points contributes to exactly one $U_{r_i}(A)$, so
\[
\binom{n}{2} = \sum_{i=1}^m U_{r_i}(A).
\]
By Lemma 89.1, each $U_{r_i}(A)\le \tfrac{\sqrt2}{2}n^{3/2}+\tfrac14 n$. Hence
\[
\binom{n}{2}\le m\Bigl(\frac{\sqrt2}{2}n^{3/2}+\frac14 n\Bigr),
\]
so
\[
D(A)=m\ge \frac{n(n-1)/2}{(\sqrt2/2)n^{3/2}+n/4}
=\frac{n-1}{\sqrt2\,\sqrt n+1/2}\,\sqrt n.
\]
For $n\ge 4$, the denominator satisfies $\sqrt2\,\sqrt n+1/2\le 2\sqrt2\,\sqrt n$, giving $m\ge \tfrac1{2\sqrt2}\sqrt n$.
\hfill $\square$

\smallskip
\noindent\textbf{FAST REALITY CHECK (integer grids).}
For an $N\times N$ integer grid ($n=N^2$ points), the distinct squared distances are exactly the distinct values of $a^2+b^2$ with $0\le a,b\le N-1$ not both zero.
A quick computation gave:
\begin{verbatim}
N=2  n=4     D=2
N=3  n=9     D=5
N=4  n=16    D=9
N=5  n=25    D=14
N=10 n=100   D=50
N=20 n=400   D=179
N=50 n=2500  D=992
N=100 n=10000 D=3663
\end{verbatim}
These values are consistent with the heuristic $D\approx n/\sqrt{\log n}$ for the grid, but no asymptotic estimate is proved here.

\medskip
\noindent\textbf{VERIFICATION.}
\begin{itemize}
\item \emph{Circle intersection:} two radius-$r$ circles are either disjoint, tangent, or meet in exactly two points; never three.
\item \emph{$K_{2,3}$-free implication:} a $K_{2,3}$ in the incidence graph would force two circles to share three points, contradicting the previous fact.
\item \emph{Algebra in Claim 2:} checked the quadratic inequality and the step $n(n-1)\le n^2$.
\item \emph{Distinct-distance pigeonhole:} the decomposition $\binom{n}{2}=\sum U_{r_i}$ holds because each pair has exactly one distance.
\end{itemize}

\medskip
\noindent\textbf{FINAL.} \textbf{UNRESOLVED}
\begin{enumerate}
\item[(i)] \textbf{Strongest proved partial result.}
For every $A\subset\mathbb{R}^2$ with $|A|=n$, one has $D(A)=\Omega(\sqrt n)$ (Lemma 89.2), via the concrete bound $U_r(A)=O(n^{3/2})$ for pairs at any fixed distance $r$ (Lemma 89.1).
\item[(ii)] \textbf{First gap (crisp).}
Improve the exponent from $1/2$ to $1$ up to logarithmic factors: prove that $D(A)\ge c\,n/\sqrt{\log n}$ for all sufficiently large $n$.
\item[(iii)] \textbf{Top 3 next moves.}
\begin{itemize}
\item Strengthen Lemma 89.1 to an upper bound $U_r(A)=O(n^{4/3})$ (or better) using incidence geometry beyond the $K_{2,3}$ extremal graph argument.
\item Prove a \emph{multi-distance energy inequality}: bound $\sum_r U_r(A)^2$ (or related additive energy) to control distance multiplicities more sharply.
\item Compute/experiment with candidate near-extremal sets (integer grids, ``perturbed'' grids, circular arcs) to identify structural features that force many repeated distances.
\end{itemize}
\item[(iv)] \textbf{Likely structure of a minimal counterexample.}
A hypothetical set $A$ with unusually small $D(A)$ must have many repeated distances, i.e. large $U_r(A)$ for many radii $r$. This suggests a highly structured configuration with many integer-like constraints (e.g. many points in approximate lattice position), rather than a generic or random configuration where almost all distances are distinct.
\end{enumerate}



% END REUSED SOURCE: attacks/erdos/gpt_pro_5.2/89.tex

\subsection*{COMPLETION ESTIMATE}
The primary source supplies no extractable numerical completion estimate. Its mathematical progress is not numerically assessed here.

Newly verified progress in this collection pass: $0\%$. This measures new verification work, not the amount of mathematics achieved in the reused text.
